%! TEX root = main.tex

% 数学相关宏包
\usepackage{amsmath, amssymb, amsthm}
\usepackage{amsfonts}
\usepackage{mathtools}
\usepackage{geometry}
\usepackage{graphicx}
\usepackage[bookmarks = false]{hyperref}
\usepackage{xcolor}
\usepackage{enumitem}
\usepackage[scheme=plain]{ctex}
\usepackage{booktabs}   % 提供 \toprule, \midrule, \bottomrule
\usepackage{caption}    % 改善图表标题间距（可选但推荐）
\usepackage{siunitx}    % 用于对齐数字（如需数值列）
\usepackage{makecell}
\usepackage{ctex}
\makeatletter
\renewcommand{\maketitle}{
    \begin{center}
        \vspace*{-1cm}  % 调整这个值来控制顶部间距
        {\LARGE\bfseries \@title \par}
        \vspace{0.2cm}  % 标题与作者间距
        {\large \@author }
%        \vspace{-0cm}  % 作者与日期间距
        {\normalsize \@date \par}
        \vspace{1cm}    % 日期与正文间距
    \end{center}
}
\makeatother

% 页面设置
\geometry{left=1.0cm,right=1.0cm,top=1.0cm,bottom=1.0cm}
% 超链接设置
\hypersetup{
    colorlinks=true,
    linkcolor=blue,
    urlcolor=blue,
    citecolor=red
}
% 定理环境定义
\newtheorem{theorem}{Theorem}[section]
\newtheorem{lemma}{Lemma}[theorem]
\newtheorem{corollary}{Corollary}[theorem]
\newtheorem{proposition}{Proposition}[theorem]
\newtheorem{definition}{Definition}[theorem]
\newtheorem{example}{Example}[theorem]
\newtheorem{remark}{Remark}[theorem]

\usepackage[noabbrev,nameinlink]{cleveref}
% 设置cleveref的标签名称
\crefname{theorem}{Theorem}{Theorems}
\Crefname{theorem}{Theorem}{Theorems}
\crefname{lemma}{Lemma}{Lemmas}
\Crefname{lemma}{Lemma}{Lemmas}
\crefname{corollary}{Corollary}{Corollaries}
\Crefname{corollary}{Corollary}{Corollaries}
\crefname{proposition}{Proposition}{Propositions}
\Crefname{proposition}{Proposition}{Propositions}
\crefname{definition}{Definition}{Definitions}
\Crefname{definition}{Definition}{Definitions}
\crefname{example}{Example}{Examples}
\Crefname{example}{Example}{Examples}
\crefname{remark}{Remark}{Remarks}
\Crefname{remark}{Remark}{Remarks}
% 数学符号定义
\newcommand{\R}{\mathbb{R}}
\newcommand{\N}{\mathbb{N}}
\newcommand{\Z}{\mathbb{Z}}
\newcommand{\Q}{\mathbb{Q}}
\newcommand{\C}{\mathbb{C}}
\newcommand{\prl}{pr_{\lambda}}
\newcommand{\pr}{pr}
\newcommand{\prt}{pr_{\theta}}
\newcommand{\fracc}[1]{\frac{#1}{2}}
\newcommand{\Filt}{\operatorname{Filt}}
\newcommand{\indentedline}[1]{%
  \par\noindent\hspace*{#1em}
}
% 自定义 enumerate 样式：使用 (1), (1.1), (1.1.1) 等格式
% 第1级：(1), (2), ...
\setlist[enumerate,1]{label=(\arabic{enumi})}
% 第2级：(1.1), (1.2), ...
\setlist[enumerate,2]{label=(\arabic{enumi}.\arabic{enumii})}
% 第3级：(1.1.1), (1.1.2), ...
\setlist[enumerate,3]{label=(\arabic{enumi}.\arabic{enumii}.\arabic{enumiii})}
% 第4级（如有需要）
\setlist[enumerate,4]{label=(\arabic{enumi}.\arabic{enumii}.\arabic{enumiii}.\arabic{enumiv})}
% 文档信息
\title{Results  }
\author{Wei}
\date{\today}

% 临时抑制 Overfull/Underfull 警告
% TEMP: suppress all overfull warnings
\hfuzz=\maxdimen
\vfuzz=\maxdimen
\hbadness=10000
\vbadness=10000


